% ---------------------------------------------------------------------------
% Figure blocks for Section IV, 2026-09-15 (lab meeting): Fig. 1 is the
% page-1 teaser (six bar panels, all 14 configurations), Fig. 3 the
% single-column trade-off grid, Fig. 4 the single-column depth and repeat
% grid, Fig. 5 the consistency chart. Files live in figure/ on Overleaf.
% Placement: the teaser block after the second paragraph of the
% Introduction (so it lands in the right column of page 1), the trade-off
% block before \subsection{Setup and Protocol}, the depth and repeat block
% after \input{tables/libero}, the consistency block before
% \subsection{Results and Discussion}. The overview block (Fig. 2, full
% width) sits right after \section{Bag of Tricks}.
% ---------------------------------------------------------------------------

\begin{figure}[!t]
  \centering
  \includegraphics[width=\columnwidth]{figure/teaser_bridge_bars.pdf}
  \caption{Success of all 14 configurations on the six SimplerEnv WidowX
  backbones, the original policy counted as one of six tricks. Within a
  trick the bars follow the order of its settings in Section~IV-A, keep
  ratio 0.2 then 0.5, two then four steps, one, two then four layers,
  strict, moderate then aggressive, and motion-entropy, task-aware then
  conservative-adaptive, darkest first. The dotted line marks the
  original, and a 0 stands where a setting succeeded on no episode. No trick raises success by more
  than 10 points on any backbone, while action repeat, depth pruning and
  foveation lower it by 10 to 85 points on several.}
  \label{fig:teaser}
\end{figure}

\begin{figure*}[t]
  \centering
  \includegraphics[width=0.85\textwidth]{figure/overview_A3.pdf}
  \caption{Where each trick enters the control loop of a frozen VLA
  policy. Candidates (top) act on the model call, the visual tokens, and the
  decoder depth, each where a signal admits it. Controls (bottom) use no
  signal. The dots and the diamond mark the points of effect.}
  \label{fig:overview}
\end{figure*}

\begin{figure}[t]
  \centering
  \includegraphics[width=\columnwidth]{figure/tradeoff_bridge_col.pdf}
  \caption{Every setting of every trick on the six SimplerEnv WidowX
  backbones, placed by its speedup per environment step and its success
  change from the original policy (the star). The shaded corner is faster beyond the latency floor and no worse than
  the original. CronusVLA's moderate and aggressive reuse settings and
  UniVLA's task-aware fusion setting are omitted because their latency is
  not read (Section~IV-A).}
  \label{fig:tradeoff-widowx}
\end{figure}

\begin{figure}[t]
  \centering
  \includegraphics[width=\columnwidth]{figure/ablation_depth_repeat_col.pdf}
  \caption{Depth pruning (a, b) and action repeat (c, d) on every
  backbone and environment, success change (a, c) and speedup (b, d)
  against the number of removed decoder layers or the number of
  times each action is repeated. Speedup is the original's time per step
  divided by the trick's,
  and LIBERO suites are pooled per backbone, so 13 lines stand for the 22
  pairs. SmolVLA's speedup at two and four layers is omitted because
  those runs used a second software stack (Section~IV-A).}
  \label{fig:ablation-depth}
\end{figure}

\begin{figure}[t]
  \centering
  \includegraphics[width=\columnwidth]{figure/cross_setting_consistency_col.pdf}
  \caption{For each setting, how many of the 13 combinations of
  Fig.~\ref{fig:ablation-depth} have lower, unchanged or higher success
  than the original, counting any nonzero change.}
  \label{fig:consistency}
\end{figure}
